% =====================================================================
% Unified Sequence Attribute Table across SatSOT, SV248S, and OOTB
% Requires: \usepackage{booktabs}, \usepackage{array}, \usepackage{makecell}
% Optional (recommended for two-column layouts): \usepackage{adjustbox}
% =====================================================================

\begin{table*}[t]
\centering
\caption{Unified sequence-attribute taxonomy across SatSOT, SV248S, and OOTB.
Rows list the unified attributes; the last three columns indicate the
corresponding per-dataset label (``---'' if the attribute is not annotated
in that dataset). Dataset-specific abbreviations appearing in the last
three columns are:
\textbf{STO} = Short-Term Occlusion,
\textbf{LTO} = Long-Term Occlusion,
\textbf{BCL} = Background Cluster ($\geq 10$ frames with the INV flag),
\textbf{IPR} = In-Plane Rotation,
\textbf{DS} = Dense Similarity,
\textbf{PO} = Partial Occlusion,
\textbf{FO} = Full Occlusion,
\textbf{SA} = Similar Appearance,
\textbf{OON} = Out-of-Normal (bounding-box aspect ratio outside $[0.3, 3]$).}
\label{tab:unified_attributes}
\small
\setlength{\tabcolsep}{4pt}
\renewcommand{\arraystretch}{1.15}
\resizebox{\textwidth}{!}{%
\begin{tabular}{llp{7.4cm}ccc}
\toprule
\textbf{Abbr.} & \textbf{Full Name} & \textbf{Definition} & \textbf{SatSOT} & \textbf{SV248S} & \textbf{OOTB} \\
\midrule
BC   & Background Clutter         & Background has similar appearance (texture/color) to the target & BC  & --- & BC  \\
BCH  & Background Change          & Background has noticeable changes in color or texture along the sequence & --- & BCH & --- \\
IV   & Illumination Variation     & Illumination around the target changes significantly & IV  & IV  & IV  \\
LQ   & Low Quality                & Image quality is low; target is hard to distinguish & LQ  & --- & --- \\
LT   & Less Textures              & Target has poor texture information, causing discrimination difficulty & --- & --- & LT  \\
MB   & Motion Blur                & Target region is blurred due to object or platform motion & --- & --- & MB  \\
ND   & Natural Disturbance        & Target appearance affected by smog, sand, or clouds & --- & ND  & --- \\
ROT  & Rotation (in-plane)        & Target rotates in the image plane ($\geq 30^{\circ}$ for SV248S) & ROT & IPR & IPR \\
POC  & Partial Occlusion$^{1}$    & Target is partially occluded & POC & STO$^{1}$ & PO  \\
FOC  & Full Occlusion$^{1}$       & Target is fully occluded    & FOC & LTO$^{1}$ & FO  \\
CO   & Continuous Occlusion       & Occlusion events occur twice or more in a sequence & --- & CO  & --- \\
INV  & Invisible$^{2}$            & Target disappears without occluder (too similar to surroundings) & --- & BCL$^{2}$ & --- \\
TO   & Tiny Object                & At least one GT bbox has fewer than 25 pixels & TO  & --- & --- \\
SOB  & Similar Object             & Nearby objects of similar shape/type/appearance & SOB & DS  & SA  \\
IM   & Isotropic Motion           & Nearby objects move with similar magnitude and direction & --- & --- & IM  \\
AM   & Anisotropic Motion         & Nearby objects move with similar magnitude but opposite direction & --- & --- & AM  \\
SM   & Slow Motion                & Target moves slowly ($<2.2$ pps in SV248S) & --- & SM  & --- \\
BJT  & Background Jitter          & Background jitter from satellite camera shaking & BJT & --- & --- \\
ARC  & Aspect Ratio Change        & Aspect ratio of bbox changes beyond a threshold (SatSOT: $[0.5, 2]$; OOTB OON: $[0.3, 3]$) & ARC & --- & OON \\
DEF  & Deformation                & Non-rigid object deformation & DEF & --- & DEF \\
\bottomrule
\end{tabular}
}
\vspace{2pt}
\begin{flushleft}\footnotesize
$^{1}$ SV248S defines STO / LTO \emph{temporally} ($\leq 50$ vs. $>50$ consecutive OCC frames),
whereas SatSOT and OOTB define partial/full occlusion \emph{spatially}. The mapping
POC $\approx$ STO and FOC $\approx$ LTO is therefore an approximation.\\
$^{2}$ SV248S BCL is derived from $\geq 10$ frames carrying the INV (Invisible) flag;
it is semantically closest to ``target indistinguishable from background'' and is kept
as a separate unified label rather than merged into BC or FOC.
\end{flushleft}
\end{table*}
